\documentclass[11pt,a4paper]{article}
\usepackage[margin=27mm]{geometry}
\usepackage[T1]{fontenc}
\usepackage{lmodern,microtype}
\usepackage{amsmath,amssymb,amsthm,mathtools}
\usepackage{booktabs,longtable,array}
\usepackage{xcolor,xurl,hyperref,enumitem,fancyhdr}
\hypersetup{colorlinks=true,linkcolor=blue!45!black,urlcolor=blue!55!black,
 citecolor=blue!45!black,pdftitle={The Basel identity in intuitionistic primitive recursive arithmetic},
 pdfauthor={Prepared for Vladimir Reshetnikov}}
\setlength{\emergencystretch}{2em}
\setlist{itemsep=3pt,topsep=5pt}
\pagestyle{fancy}\fancyhf{}
\fancyhead[L]{\small The Basel identity in iPRA}
\fancyhead[R]{\small A finite rational proof}
\fancyfoot[C]{\thepage}\setlength{\headheight}{14pt}
\newtheorem{theorem}{Theorem}[section]
\newtheorem{lemma}[theorem]{Lemma}
\newtheorem{proposition}[theorem]{Proposition}
\newtheorem{corollary}[theorem]{Corollary}
\theoremstyle{remark}\newtheorem{remark}[theorem]{Remark}
\newcommand{\N}{\mathbb N}
\newcommand{\ipra}{\mathsf{iPRA}}
\newcommand{\abs}[1]{\left|#1\right|}
\newcommand{\cut}{\mathcal N}
\newcommand{\dm}{\mathbin{\dot{-}}}
\DeclareMathOperator{\Prf}{Prf}

\title{\textbf{The Basel identity in intuitionistic\\
primitive recursive arithmetic}\\[5pt]
\large A complete finite rational argument and a Lean formalization blueprint}
\author{Prepared for Vladimir Reshetnikov}
\date{November 14, 2026}
\begin{document}
\maketitle
\begin{abstract}
We prove Euler's identity $\sum_{j\ge1}j^{-2}=\pi^2/6$ in its explicit
first-order arithmetical form. Put
$A_N=\sum_{j<N}(-1)^j/(2j+1)$ and $S_n=\sum_{j=1}^n j^{-2}$.
The central result is the finite, uniform inequality
\[
 n\ge2^r\quad\Longrightarrow\quad
 \abs{6S_n-16A_{2^r}^{\,2}}\le\frac{8r+20}{2^r}.
\]
Consequently, with $\cut(k)=2^{2k+7}$ and $P_k=4A_{\cut(k)}$,
$n\ge\cut(k)$ implies $\abs{6S_n-P_k^2}\le2^{-k}$.
All sums are finite, all arithmetic is rational, and all inductions can
be performed on quantifier-free formulas. The proof compares two
decompositions of the same finite square and bounds both boundary
terms. A separate finite Riemann-sum calculation identifies the
chosen $\pi$ code with the usual arctangent definition. We supply
natural-number encodings, an induction audit, exact computational
checks, and a concrete path to an object-theory derivation using
ConciseIPRA. The mathematical argument is complete. The Lean files
are an explicitly unfinished, uncompiled formalization blueprint,
not a claimed machine-checked iPRA proof.
\end{abstract}
\setcounter{tocdepth}{1}
\begingroup\small\tableofcontents\endgroup
\newpage

\section{The precise theorem and its interpretation}\label{sec:statement}

The first-order language of $\ipra$ has no variables for real numbers,
infinite sequences, or sets. Its internal Basel statement must concern
specified primitive recursive programs returning rational
approximations. This is a change of representation, not a weakening to
numerical evidence.

For $N,n\in\N$, define the finite rational sums
\begin{align}
 A_N&=\sum_{j=0}^{N-1}\frac{(-1)^j}{2j+1},&
 H_N&=\sum_{j=0}^{N-1}\frac1{2j+1},\label{eq:AH}\\
 O_N&=\sum_{j=0}^{N-1}\frac1{(2j+1)^2},&
 S_n&=\sum_{j=1}^{n}\frac1{j^2}.\label{eq:OS}
\end{align}
An empty sum is zero. The following result is proved entirely by
finite rational arithmetic.

\begin{theorem}[Uniform finite Basel estimate]\label{thm:uniform}
For all natural numbers $n\ge N\ge1$,
\begin{equation}\label{eq:uniform}
 \boxed{\displaystyle\abs{6S_n-16A_N^2}\le\frac{16H_N+4}{N}.}
\end{equation}
For all $r,n\in\N$ with $n\ge2^r$,
\begin{equation}\label{eq:dyadic}
 \boxed{\displaystyle\abs{6S_n-16A_{2^r}^{\,2}}\le\frac{8r+20}{2^r}.}
\end{equation}
\end{theorem}

\begin{theorem}[An explicit first-order precision statement]\label{thm:main}
Set
\begin{equation}\label{eq:cutoff}
 \cut(k)=2^{2k+7},\qquad P_k=4A_{\cut(k)}.
\end{equation}
Then $\ipra$ proves the universal closure of
\begin{equation}\label{eq:main}
 \boxed{\displaystyle n\ge\cut(k)\ \longrightarrow\
 \abs{6S_n-P_k^2}\le2^{-k}.}
\end{equation}
Rational notation abbreviates a quantifier-free natural-number
formula specified in Section~\ref{sec:encoding}.
\end{theorem}

The ordinary interpretation of $P_k$ converges to $\pi$. We prove its
finite Cauchy estimate and its agreement with an independent rational
Riemann-sum code in Section~\ref{sec:pi}. Thus \eqref{eq:main} is a
constructive arithmetical form of the exact Basel identity, with a
primitive recursive precision witness already supplied.

Squaring Gregory's series belongs to a historical family of proofs.
Campbell and Levrie discuss that family, including Euler's 1742
correspondence with Goldbach~\cite{CampbellLevrie}. The point here is
the explicit finite error control and proof-theoretic accounting.
No claim of historical novelty is needed.

\section{Elementary finite arithmetic}\label{sec:elementary}

Through Section~\ref{sec:modulus}, every quantity is a rational number
computed by a finite expression or bounded loop. Absolute value is
decidable rational absolute value. All denominators used to multiply
inequalities are positive.

\begin{lemma}[Decreasing alternating terms]\label{lem:decreasing}
If $w_0,\ldots,w_{m-1}$ are nonnegative and decreasing, then
\[
 0\le w_0-w_1+w_2-\cdots+(-1)^{m-1}w_{m-1}\le w_0
 \quad(m\ge1).
\]
An alternating sum over any finite interval therefore has absolute
value at most its first, largest term.
\end{lemma}
\begin{proof}
Process the list backwards. Set $R_0=0$ and
$R_{j+1}=w_{m-1-j}-R_j$ for $0\le j<m$. The first step gives
$0\le R_1\le w_{m-1}$. If $0\le R_j\le w_{m-j}$ and $j<m$,
then $w_{m-j}\le w_{m-1-j}$ and consequently
$0\le R_{j+1}\le w_{m-1-j}$. Finally $R_m$ is the displayed
alternating sum. Multiplication by the sign of the first term
handles an interval beginning at an odd index.
\end{proof}

\begin{lemma}[Increasing alternating terms]\label{lem:increasing}
If $0\le b_1\le\cdots\le b_m$, then
\[
 \abs{\sum_{d=1}^m(-1)^{d-1}b_d}\le b_m.
\]
\end{lemma}
\begin{proof}
Set $U_0=0$ and $U_d=b_d-U_{d-1}$. Induction gives
$0\le U_d\le b_d$, since $0\le U_{d-1}\le b_{d-1}\le b_d$.
The recurrence also gives
$\sum_{j=1}^d(-1)^{j-1}b_j=(-1)^{d-1}U_d$.
Both invariants are equalities or inequalities of finite rational
terms.
\end{proof}

These are finite lemmas, not convergence theorems. In their
applications the lists are explicit functions of natural-number
parameters, not function variables in the language of $\ipra$.

The partial-fraction calculations needed below are
\begin{align}
 \frac1{ab}&=\frac{1/a+1/b}{a+b}\quad(a,b>0),\label{eq:pfplus}\\
 \frac2{(2i+1)(2i+2d+1)}
 &=\frac1d\left(\frac1{2i+1}-\frac1{2i+2d+1}\right)
 \quad(d\ge1).\label{eq:pfminus}
\end{align}
They follow by cross-multiplication. Every subsequent reindexing is
finite and can be checked by adding a row, column, or diagonal.

\section{The first decomposition: sum diagonals}\label{sec:triangle}

Fix $N\ge1$ and put
\begin{equation}\label{eq:TN}
 T_N=\sum_{d=1}^{N-1}\frac{(-1)^{d-1}H_d}{d}.
\end{equation}
\begin{lemma}[A finite convolution identity]\label{lem:diagonal}
For $d\ge1$,
\begin{equation}\label{eq:diagonal}
 \sum_{i=0}^{d-1}\frac1{(2i+1)(2d-2i-1)}=\frac{H_d}{d}.
\end{equation}
Consequently
\begin{equation}\label{eq:triangle}
 T_N=\sum_{\substack{0\le i,j<N\\i+j\le N-2}}
       \frac{(-1)^{i+j}}{(2i+1)(2j+1)}.
\end{equation}
For $N=1$, the last sum is empty.
\end{lemma}
\begin{proof}
The two denominators in \eqref{eq:diagonal} add to $2d$, so its
summand is
\[
 \frac1{2d}\left(\frac1{2i+1}+\frac1{2d-2i-1}\right).
\]
Both sums in parentheses enumerate the same $d$ odd reciprocals,
one in reverse order. This proves \eqref{eq:diagonal}.
In \eqref{eq:triangle}, collect the terms with $i+j=d-1$.
Their sign is $(-1)^{d-1}$ and their unsigned sum is $H_d/d$.
The permitted values are exactly $d=1,\ldots,N-1$.
\end{proof}

The square $A_N^2$ is the sum over the full $N$ by $N$ square.
Write $D_N=A_N^2-T_N$ for its complementary triangular part.

\begin{lemma}[The omitted triangle]\label{lem:D}
For every $N\ge1$,
\begin{equation}\label{eq:D}
 \abs{D_N}\le\frac{H_N}{N}.
\end{equation}
\end{lemma}
\begin{proof}
In row $i$, the omitted indices are $j=N-1-i,\ldots,N-1$.
Hence
\[
 D_N=\sum_{i=0}^{N-1}\frac{(-1)^i}{2i+1}
       \sum_{j=N-1-i}^{N-1}\frac{(-1)^j}{2j+1}.
\]
Lemma~\ref{lem:decreasing} bounds the inner absolute value by
$1/(2N-2i-1)$. The finite triangle inequality and
\eqref{eq:diagonal}, with $d=N$, give
\[
 \abs{D_N}\le
 \sum_{i=0}^{N-1}\frac1{(2i+1)(2N-2i-1)}=\frac{H_N}{N}.
\]
\end{proof}

\section{The second decomposition: difference diagonals}\label{sec:difference}

Split the same square into its main diagonal and pairs above and below
it. If $j=i+d$, then $(-1)^{i+j}=(-1)^d$, so
\begin{equation}\label{eq:square}
 A_N^2=O_N+
 2\sum_{d=1}^{N-1}(-1)^d
 \sum_{i=0}^{N-d-1}\frac1{(2i+1)(2i+2d+1)}.
\end{equation}
For fixed $d$, equation \eqref{eq:pfminus} yields
\begin{align}
 2\sum_{i=0}^{N-d-1}\frac1{(2i+1)(2i+2d+1)}
 &=\frac{H_{N-d}-(H_N-H_d)}d\notag\\
 &=\frac{H_d}d-\frac{H_N-H_{N-d}}d.\label{eq:difference}
\end{align}
Define
\begin{equation}\label{eq:bE}
 b_d=\frac{H_N-H_{N-d}}d\quad(1\le d<N),\qquad
 E_N=\sum_{d=1}^{N-1}(-1)^{d-1}b_d.
\end{equation}
Then \eqref{eq:TN}--\eqref{eq:difference} give the exact equality
\begin{equation}\label{eq:second}
 \boxed{A_N^2=O_N-T_N+E_N.}
\end{equation}

\begin{lemma}[Increasing tail averages]\label{lem:average}
For fixed $N\ge2$, the $b_1,\ldots,b_{N-1}$ are nonnegative and
increasing.
\end{lemma}
\begin{proof}
Let $C_d=H_N-H_{N-d}$. For $1\le d<N-1$, set
$t=1/(2(N-d-1)+1)$. Then $C_{d+1}=C_d+t$.
Each of the $d$ summands in $C_d$ has a larger denominator than
$t$ does, so $C_d\le dt$. Therefore
\[
 b_{d+1}-b_d
 =\frac{C_d+t}{d+1}-\frac{C_d}{d}
 =\frac{dt-C_d}{d(d+1)}\ge0.
\]
Nonnegativity follows from the defining sum.
\end{proof}

\begin{lemma}[The second boundary term]\label{lem:E}
For every $N\ge1$,
\begin{equation}\label{eq:E}
 \abs{E_N}\le\frac{H_N}{N}.
\end{equation}
\end{lemma}
\begin{proof}
For $N=1$, the sum is empty. For $N\ge2$, the preceding lemma and
Lemma~\ref{lem:increasing} imply
\[
 \abs{E_N}\le b_{N-1}=\frac{H_N-1}{N-1}.
\]
Every summand of $H_N$ is at most one, so $H_N\le N$.
After multiplication by $N(N-1)>0$, the inequality
$(H_N-1)/(N-1)\le H_N/N$ is exactly $H_N\le N$.
\end{proof}

\begin{proposition}[The crucial finite comparison]\label{prop:core}
For every $N\ge1$,
\begin{equation}\label{eq:core}
 \boxed{\displaystyle\abs{2A_N^2-O_N}\le\frac{2H_N}{N}.}
\end{equation}
\end{proposition}
\begin{proof}
Since $D_N=A_N^2-T_N$, equation \eqref{eq:second} implies
$2A_N^2-O_N=D_N+E_N$. Apply \eqref{eq:D} and \eqref{eq:E}.
\end{proof}

This step evaluates the series. It compares two finite decompositions
and retains both boundary errors. Formally squaring an infinite
conditionally convergent series would not justify dropping those
errors.

\section{From odd squares to all squares}\label{sec:allsquares}

\begin{lemma}[Finite square tails]\label{lem:tail}
For $m\ge N\ge1$,
\begin{equation}\label{eq:tail}
 0\le S_m-S_N\le\frac1N-\frac1m\le\frac1N,
\end{equation}
and
\begin{equation}\label{eq:annulus}
 \frac1{4N}\le S_{2N}-S_N\le\frac1{2N}.
\end{equation}
\end{lemma}
\begin{proof}
For $j\ge2$,
$1/j^2\le1/(j(j-1))=1/(j-1)-1/j$.
Sum from $N+1$ to $m$ to get \eqref{eq:tail}.
Its case $m=2N$ is the upper bound in \eqref{eq:annulus}.
For the lower bound there are $N$ terms, each at least
$1/(2N)^2$.
\end{proof}

Separating the even and odd denominators gives
\begin{equation}\label{eq:oddeven}
 O_N=S_{2N}-\frac14S_N.
\end{equation}

\begin{lemma}[Uniform comparison]\label{lem:compare}
For $n\ge N\ge1$,
\begin{equation}\label{eq:compare}
 \abs{O_N-\frac34S_n}\le\frac1{2N}.
\end{equation}
\end{lemma}
\begin{proof}
Put $R=S_{2N}-S_N$ and $V=S_n-S_N$. The previous lemma says
\[
 \frac1{4N}\le R\le\frac1{2N},\qquad 0\le V\le\frac1N.
\]
By \eqref{eq:oddeven}, $O_N-\frac34S_n=R-\frac34V$.
Its upper bound is $1/(2N)$ and its lower bound is
$1/(4N)-3/(4N)=-1/(2N)$.
\end{proof}

\begin{proof}[Proof of \eqref{eq:uniform}]
The core comparison and \eqref{eq:compare} imply
\[
 \abs{\frac34S_n-2A_N^2}
 \le\abs{\frac34S_n-O_N}+\abs{O_N-2A_N^2}
 \le\frac1{2N}+\frac{2H_N}{N}.
\]
Multiply by eight.
\end{proof}

\section{The primitive recursive precision witness}\label{sec:modulus}

\begin{lemma}[Dyadic harmonic bound]\label{lem:harmonic}
For every $r\in\N$,
\begin{equation}\label{eq:harmonic}
 H_{2^r}\le1+\frac r2.
\end{equation}
\end{lemma}
\begin{proof}
$H_1=1$. Moreover,
\[
 H_{2^{r+1}}-H_{2^r}
 =\sum_{j=2^r}^{2^{r+1}-1}\frac1{2j+1}
 \le2^r\cdot2^{-r-1}=\frac12.
\]
Induction proves the claim; after clearing denominators the invariant
is $2H_{2^r}\le r+2$.
\end{proof}

Substitution in \eqref{eq:uniform} proves \eqref{eq:dyadic}.
Next,
\begin{equation}\label{eq:growth}
 16k+76\le128\cdot2^k\quad(k\in\N).
\end{equation}
The base case is $76\le128$. At the successor,
$16(k+1)+76\le2(16k+76)\le128\cdot2^{k+1}$.
This is quantifier-free induction on a natural-number inequality.

Take $r=2k+7$ in \eqref{eq:dyadic}. Its numerator is $16k+76$,
so \eqref{eq:growth} gives
\[
 \frac{8r+20}{2^r}
 \le\frac{128\cdot2^k}{2^{2k+7}}=2^{-k}.
\]
This proves the rational assertion \eqref{eq:main}. Section
\ref{sec:encoding} spells out its availability in $\ipra$.

\section{What the \texorpdfstring{$\pi$}{pi} code means}\label{sec:pi}

\subsection{Finite Cauchy estimates}

Lemma~\ref{lem:decreasing} gives, for $M\ge N\ge1$,
\begin{equation}\label{eq:gregorytail}
 0\le A_N\le1,\qquad
 \abs{A_M-A_N}\le\frac1{2N+1}.
\end{equation}
Thus $P_k=4A_{\cut(k)}$ has an explicit Cauchy modulus: for
$\ell\ge k$,
\[
 \abs{P_\ell-P_k}\le\frac4{2\cut(k)+1}<\frac2{\cut(k)}.
\]
Likewise
$\abs{S_{\cut(\ell)}-S_{\cut(k)}}\le1/\cut(k)$
by \eqref{eq:tail}. These are statements about specified arithmetic
programs. They do not assert that every Cauchy sequence has a limit.

In the usual real interpretation, $A_N\to\pi/4$. The next calculation
identifies the code with a rational Riemann-sum definition, without
using calculus inside the arithmetic proof.

\subsection{A finite arctangent bridge}\label{sec:riemann}

For $M\ge1$ define
\begin{equation}\label{eq:QM}
 Q_M=\frac1M\sum_{i=0}^{M-1}\frac1{1+(i/M)^2}
     =M\sum_{i=0}^{M-1}\frac1{M^2+i^2}.
\end{equation}
This rational program is independent of the Basel sums.

\begin{lemma}[Power-sum quadrature]\label{lem:moment}
For $p\in\N$, $M\ge1$, and $0^0=1$,
\[
 0\le\frac1{p+1}
       -\frac1M\sum_{i=0}^{M-1}(i/M)^p\le\frac1M.
\]
\end{lemma}
\begin{proof}
For $p=0$ there is equality. For $p\ge1$, the finite factorization
of a difference of powers gives
\[
 (p+1)i^p\le(i+1)^{p+1}-i^{p+1}\le(p+1)(i+1)^p.
\]
Each of the $p+1$ terms in that factorization is between $i^p$
and $(i+1)^p$. Sum over $i=0,\ldots,M-1$ and divide by
$(p+1)M^{p+1}$. The two endpoint sums differ by $1/M$,
which proves the bound.
\end{proof}

\begin{proposition}[Rational identification]\label{prop:quadrature}
For $N,M\ge1$,
\begin{equation}\label{eq:quadrature}
 \abs{Q_M-A_N}\le\frac NM+\frac1{2N+1}.
\end{equation}
\end{proposition}
\begin{proof}
For rational $x\in[0,1]$, multiplication by $1+x^2$ proves
the finite geometric identity
\begin{equation}\label{eq:geometric}
 \frac1{1+x^2}
 =\sum_{j=0}^{N-1}(-1)^j x^{2j}
  +(-1)^N\frac{x^{2N}}{1+x^2}.
\end{equation}
Substitute $x=i/M$ and average over $i<M$.
Replacing each power average by $1/(2j+1)$ costs at most $1/M$
by Lemma~\ref{lem:moment}, for a total cost at most $N/M$.
The remainder average has absolute value at most
\[
 \frac1M\sum_{i=0}^{M-1}(i/M)^{2N}\le\frac1{2N+1}.
\]
Add the two bounds.
\end{proof}

The $Q_M$ also have their own finite Cauchy estimate. The rational
function $f(x)=1/(1+x^2)$ is decreasing on $[0,1]$ by
cross-multiplication. Its lower and upper rectangular sums differ
by $(f(0)-f(1))/M=1/(2M)$. Refinement into $ML$ pieces keeps a
rectangular sum between the original lower and upper sums.
Grouping the $ML$ summands into $M$ blocks and then into $L$ blocks
therefore gives
\begin{equation}\label{eq:Qmodulus}
 \abs{Q_M-Q_L}\le\frac1{2M}+\frac1{2L}.
\end{equation}
This is a finite comparison of sums, not an appeal to existence of
an integral.

With $M=N^2$, \eqref{eq:quadrature} yields
$\abs{4Q_{N^2}-4A_N}\le6/N$. Thus the two explicit arithmetic
codes name the same constant, with an explicit comparison modulus.

\begin{corollary}[An independent code works as well]\label{cor:riemann}
Set $\widetilde P_k=4Q_{\cut(k)^2}$. Then $\ipra$ proves
\[
 n\ge\cut(k)\quad\Longrightarrow\quad
 \abs{6S_n-\widetilde P_k^{\,2}}\le2^{-k}.
\]
\end{corollary}
\begin{proof}
Since $0\le Q_M,A_N\le1$,
\[
 \abs{16Q_{N^2}^{\,2}-16A_N^2}
 \le32\left(\frac1N+\frac1{2N+1}\right)\le\frac{48}{N}.
\]
For $N=2^r$, combine this with \eqref{eq:dyadic}, obtaining
$(8r+68)/2^r$. At $r=2k+7$ the numerator is $16k+124$.
The inequality $16k+124\le128\cdot2^k$ follows by the same
quantifier-free induction as \eqref{eq:growth}.
\end{proof}

\subsection{The conventional geometric interpretation}

Externally, \eqref{eq:QM} is the left rectangular sum for
$\int_0^1(1+x^2)^{-1}\,dx=\arctan(1)=\pi/4$.
The unit-circle interpretation can also be seen from the rational
parametrization
\[
 \left(\frac{1-t^2}{1+t^2},\frac{2t}{1+t^2}\right),
 \qquad 0\le t\le1.
\]
It traverses a quarter circle and has speed $2/(1+t^2)$, giving
arc length $2\arctan(1)=\pi/2$. This identifies the standard meaning
of the constant. It is not a use of calculus inside $\ipra$:
the internal relation between the rational codes was already proved
in Proposition~\ref{prop:quadrature}.

In that usual interpretation,
$\abs{4A_N-\pi}\le4/(2N+1)$ and $0\le\pi\le4$, whence
$\abs{16A_N^2-\pi^2}\le16/N$. Consequently
\begin{equation}\label{eq:realmodulus}
 n\ge2^r\quad\Longrightarrow\quad
 \abs{6S_n-\pi^2}\le\frac{8r+36}{2^r}.
\end{equation}
The same cutoff $\cut(k)$ works, since
$16k+92\le128\cdot2^k$. In particular,
\[
 \boxed{\displaystyle \sum_{j=1}^{\infty}\frac1{j^2}=\frac{\pi^2}{6}.}
\]
The internal assertion consists of the finite estimates, so it does
not require the general completeness of any real-number system.

\section{Why the proof is available in iPRA}\label{sec:encoding}

\subsection{The theory}

We use intuitionistic first-order logic with equality; primitive
recursive function symbols and their defining equations; successor
injectivity and nonzeroness; and induction for quantifier-free
formulas with parameters.
This is the presentation in the inspected
\texttt{ConciseIPRA/PR.lean} and \texttt{Deduction.lean}, at
repository commit
\begin{center}\small
\texttt{b03925cccc76d55b9dc2a00864a5e0544b64aa7f}.
\end{center}
The source references and transfer theorem are recorded below
\cite{ProveItPR,ProveItNative}.

Only decidable integer or rational comparisons are split into cases.
Such case splits are constructively provable. There is no excluded
middle for an arbitrary formula, choice principle, or minimization
without an explicit bound.

\subsection{Rational codes}

Represent a signed rational by $(a,b,d)\in\N^3$, $d>0$, with value
$(a-b)/d$. No greatest-common-divisor reduction is needed.
For example,
\begin{align*}
 (a,b,d)+(c,e,f)&=(af+cd,\;bf+ed,\;df),\\
 (a,b,d)(c,e,f)&=(ac+be,\;ae+bc,\;df).
\end{align*}
Negation swaps the first two components. Equality and order are
natural-number tests:
\begin{align*}
 (a-b)/d=(c-e)/f&\ \Longleftrightarrow\ af+ed=bf+cd,\\
 (a-b)/d\le(c-e)/f&\ \Longleftrightarrow\ af+ed\le bf+cd.
\end{align*}
Division by a positive natural $t$ multiplies the denominator by $t$.
Absolute value uses a decidable comparison of numerator components.
Definitions can be totalized on invalid inputs; the proof uses valid
codes and proves denominator positivity along with its invariants.

Addition, multiplication, truncated subtraction, order tests, parity,
powers, factorials, bounded sums, bounded products, and bounded
division are primitive recursive. Their required elementary laws
are proved using defining equations and quantifier-free induction.
For example, quotient and remainder can be computed by a loop of
length at most the dividend, preserving $a=qd+r$ and $r<d$ when
$d>0$. None of the definitions relies on general recursion.

\subsection{An explicit natural-number final formula}

Define a natural fraction sum by
\[
 (u,v)\oplus(a,b)=(ub+av,vb),
 \qquad(0,1)\text{ for the empty sum}.
\]
Let $(s_{\rm num}(n),s_{\rm den}(n))$ be the fold of $(1,j^2)$,
$1\le j\le n$. Let $(p_{\rm num}(k),p_{\rm den}(k))$ be the fold of
\[
 \left(8,(4j+1)(4j+3)\right),\qquad0\le j<\cut(k)/2.
\]
Each component is primitive recursive. As $\cut(k)$ is even,
pairing adjacent terms and using
$4(1/(4j+1)-1/(4j+3))=8/((4j+1)(4j+3))$ proves
\[
 S_n=\frac{s_{\rm num}(n)}{s_{\rm den}(n)},\qquad
 P_k=\frac{p_{\rm num}(k)}{p_{\rm den}(k)}.
\]
The denominators are positive by product induction.

Write $\Delta(x,y)=(x\dm y)+(y\dm x)$.
Abbreviate $u=s_{\rm num}(n)$, $v=s_{\rm den}(n)$,
$a=p_{\rm num}(k)$, $b=p_{\rm den}(k)$.
Formula \eqref{eq:main} is precisely
\begin{equation}\label{eq:natmain}
 \boxed{\displaystyle
 \cut(k)\le n\ \longrightarrow\
 2^k\Delta(6ub^2,a^2v)\le vb^2.}
\end{equation}
Its universal closure has a quantifier-free matrix and only
primitive recursive function symbols.
\texttt{code/arithmetic\_spec.py} evaluates this representation.
The Lean specification additionally uses signed rational triples
for the intermediate lemmas.

Appendix~\ref{app:factorial} gives an alternative formula with one
factorial denominator. It is the recommended final checker target in
the supplied Lean outline. Either formula represents the same
finite theorem.

\subsection{The induction audit}\label{sec:induction}

Primitive recursive computability by itself does not imply
provability of correctness in $\ipra$. Here correctness comes from
the displayed elementary inductions. The table records their
invariants.

\begin{longtable}{@{}>{\raggedright\arraybackslash}p{.23\textwidth}
>{\raggedright\arraybackslash}p{.46\textwidth}
>{\raggedright\arraybackslash}p{.24\textwidth}@{}}
\toprule
Calculation&Invariant or local identity&Fixed parameters\\
\midrule\endfirsthead
\toprule Calculation&Invariant or local identity&Fixed parameters\\
\midrule\endhead
Rational arithmetic&Cross-multiplied identities; denominator positivity&All operands\\
Finite sum algebra&Equality of the two sums after a specified number of terms&Endpoints and summand parameters\\
Decreasing alternating row&$0\le R_j\le w_{m-j}$ in the backward recurrence&Row endpoints and $N$\\
Increasing alternating list&$0\le U_d\le b_d$ and the parity identity&$N$\\
Sum diagonals&Processed triangle equals its processed diagonal sum&Outer square size\\
Difference diagonals&Processed off-diagonal pairs equal their partial-fraction sum&Outer square size\\
Increasing averages&$C_d\le dt$ from the $d$ termwise bounds&$N,d,t$\\
Square tails&Partial tail is at most the telescoping expression&Starting index $N$\\
Dyadic blocks&$2H_{2^r}\le r+2$&None\\
Precision witness&$16k+c\le128\cdot2^k$ for $c=76,92,124$&Chosen constant\\
Finite quadrature&Telescoping power identity and endpoint bounds&$p,M$\\
\bottomrule
\end{longtable}

After rational coding, each invariant is quantifier-free. A finite
sum is a primitive recursive term, not an unbounded existential
assertion. Reindexing can be implemented by a rectangular loop
with a primitive recursive $0/1$ test for membership in a triangle.
The row and diagonal identities then follow from the loop
recurrences. A bounded collection of already checked entries may,
if useful, be represented by a Boolean fold; the induction formula
is equality of the fold to one.

Keep $n,N,k$ free while performing these inductions. Do not induct
on a statement that quantifies over every later partial sum or
every precision. Prove the quantifier-free statement for the fixed
parameters first, then universally generalize. This distinction
eliminates a common hidden increase in induction strength.

Consequently every lemma is derivable in the definitional extension
of $\ipra$ by the specified primitive recursive functions.
Replacing the defined symbols by their primitive recursive syntax
gives a derivation in the original language. Positive
cross-multiplication and universal introduction give
\eqref{eq:natmain}, proving Theorem~\ref{thm:main}.
There is no appeal to the soundness or consistency of $\ipra$.

\section{A concrete Lean formalization path}\label{sec:lean}

\subsection{The certification boundary}

A Lean theorem about \texttt{Rat} or \texttt{Real} does not itself
prove derivability in $\ipra$. The two distinct tasks are:
\begin{enumerate}
\item formalize the finite identities and estimates;
\item construct an object-theory proof term
$\Prf(\mathrm{PRAAx},[],\mathrm{BaselFormula})$.
\end{enumerate}
The second certifies the weak-arithmetic claim. The first is useful
for checking mathematics, but cannot be substituted for the second.

\subsection{Inspected repository support}

The namespace is \texttt{LeanProofs.ConciseIPRA}. At the pinned
revision the relevant interfaces are:
\begin{itemize}
\item \texttt{PR.lean}: \texttt{PR.zero}, \texttt{PR.succ},
\texttt{PR.proj}, \texttt{PR.comp}, \texttt{PR.prec},
terms \texttt{PTm}, axioms \texttt{PRAAx}, and
\texttt{PRA.qfInd}.
\item \texttt{Deduction.lean}: intuitionistic natural deduction
\texttt{Prf Ax Gamma phi}.
\item \texttt{Lindenbaum.lean}: terms modulo provable equality,
with \texttt{Q.eqInd} and \texttt{Q.exact}.
\item \texttt{TeqArith.lean}, \texttt{TeqOrder.lean},
\texttt{TeqPair.lean}, \texttt{TeqRec.lean}, and
\texttt{TeqReflect.lean}: arithmetic, pairing, recursion, and
quantifier-free reflection infrastructure.
\item \texttt{NativeBiInterpretation.lean}: the existing
\texttt{pra\_iff\_native} theorem.
\end{itemize}
The root toolchain is \texttt{leanprover/lean4:v4.32.0}.
This infrastructure is useful, but is not already a rational-analysis
library containing the present finite lemmas.

The final goal is
\begin{quote}\small
\texttt{([] : List (Fm PTm 0))} $\vdash$[\texttt{PRAAx}]
\texttt{baselFormula}.
\end{quote}
The native concise-core theorem follows by
\texttt{pra\_iff\_native.mp}. No new bi-interpretation is required.

\subsection{Implementation order}

\begin{enumerate}[label=\textbf{\arabic*.},leftmargin=*]
\item \textbf{Compile the concrete programs.}
Implement powers, parity, truncated subtraction, distance,
comparison, the fraction folds, and, for the recommended target,
factorials and bounded division in \texttt{PR} syntax.
Prove defining equations and representation lemmas inside the
object theory. Merely proving a Boolean checker true in the Lean
metatheory does not supply the internal derivation.

\item \textbf{Build arithmetic support.}
Prove the required ordered semiring identities and rational
cross-multiplication lemmas by quantifier-free induction.
An alternative is to extend the existing \texttt{Q Gamma}
arithmetic and export equalities using \texttt{Q.exact}.
Because that quotient is defined from provability, this can preserve
the certification boundary. Its algebraic instances must themselves
be established.

\item \textbf{Implement alternating cancellation.}
Use the two recurrences in Section~\ref{sec:elementary}, instantiated
with the explicit odd reciprocals and tail averages. No object-level
function variables are required.

\item \textbf{Prove the two finite decompositions.}
Represent triangles by bounded double sums. Prove
\eqref{eq:diagonal} and \eqref{eq:difference} by partial fractions,
reversal, and telescoping. Establish \eqref{eq:triangle},
\eqref{eq:square}, and \eqref{eq:second} by row and diagonal
recurrences. This is the main finite combinatorial implementation
work.

\item \textbf{Assemble the inequalities.}
Follow Lemmas~\ref{lem:D}--\ref{lem:E},
Proposition~\ref{prop:core}, and
Lemmas~\ref{lem:tail}--\ref{lem:compare}.
After positivity is recorded, the remaining calculations are
finite-sum inequalities and polynomial arithmetic.

\item \textbf{Prove the modulus and close the sentence.}
Use \eqref{eq:harmonic} and \eqref{eq:growth}, instantiate
$N=2^{2k+7}$, cross-multiply, and universally generalize.
Construct the actual \texttt{PRAAx} proof term, then apply the
existing native bridge.

\item \textbf{Add the interpretation separately.}
Formalize Section~\ref{sec:riemann} if the independent integral
code is desired. An optional external mathlib proof identifies
either code with \texttt{Real.pi} and states the usual infinite
sum. That external interpretation is not a premise of the
arithmetical proof.
\end{enumerate}

This path leaves no unspecified analytic result to prove. It does
leave substantial formal work on arithmetic coding, reindexing,
and object-theory derivations.

\subsection{Status of the supplied source}

The files in the \texttt{lean} directory have distinct roles:
\begin{itemize}
\item \texttt{BaselSpec.lean} specifies the sums and remaining
finite proof obligations.
\item \texttt{ObjectTheoryTarget.lean} uses the repository interfaces
and supplies only conditional assembly and transfer.
\item \texttt{ROADMAP.md} records the exact remaining tasks.
\end{itemize}

No completed proof of the finite Basel estimate, compiled
\texttt{PR} checker, or completed \texttt{PRAAx} derivation is
claimed. The Lean files have not been compiled here: Lean and
Lake are unavailable in this environment. They contain no
\texttt{sorry} or \texttt{axiom} declarations, but that is not
evidence of completeness; remaining obligations are explicit,
uninhabited specifications or assumptions to conditional results.

\section{Reproducibility and validation}

The archive contains this source and PDF, exact rational checks,
natural-number specifications, finite-quadrature checks, the Lean
outline, provenance, and build instructions.

The main exact checks cover $N=1,\ldots,80$ and dyadic $N$ through
$4096$. They include $3240$ omitted-row checks, $6320$ diagonal
checks, $428$ uniform-bound samples, and $256$ integer-modulus
samples. Every mathematical comparison uses exact rational or
integer arithmetic. These checks catch signs, endpoint mistakes,
and transcription errors. They are not the proof of the universal
theorem; that proof is the preceding finite argument.

The repository revision and inspected paths are recorded in
\texttt{validation/provenance.json}. No remote repository changes
are part of the package.

\appendix
\section{A common-denominator natural-number sentence}\label{app:factorial}

Define
\[
 F(n)=(2n)!,\qquad
 B(n)=\sum_{j=1}^{n}\bigl(F(n)\mathbin{\mathrm{div}}j\bigr)^2,
\]
and, with $N=\cut(k)$,
\[
 C(k,n)=\sum_{j=0}^{N/2-1}
 \left[
   \bigl(F(n)\mathbin{\mathrm{div}}(4j+1)\bigr)
   \dm
   \bigl(F(n)\mathbin{\mathrm{div}}(4j+3)\bigr)
 \right].
\]
These are total natural-number functions. All displayed divisors
are positive. If $n\ge N$, every divisor is at most $2n$, hence
divides $F(n)$. The quotient by $4j+1$ is at least the quotient
by $4j+3$, so truncated subtraction is exact. Therefore
\[
 \frac{B(n)}{F(n)^2}=S_n,\qquad
 \frac{C(k,n)}{F(n)}=A_{\cut(k)}.
\]
The first identity is termwise and the second pairs adjacent
Gregory terms. Factorial divisibility follows by finite product
factorization.

Theorem~\ref{thm:main} is thus equivalently
\begin{equation}\label{eq:factorialmain}
 \boxed{\displaystyle
 \forall k\,\forall n\
 \left(\cut(k)\le n\ \longrightarrow\
 2^k\Delta(6B(n),16C(k,n)^2)\le F(n)^2\right).}
\end{equation}
Only powers, factorials, bounded division, bounded sums, truncated
subtraction, and order occur. The large intermediate integers do not
affect primitive recursiveness.

\section{A small exact example}

For $N=2$,
\[
 A_2=\frac23,\quad H_2=\frac43,\quad O_2=\frac{10}{9},
 \quad T_2=1,\quad D_2=-\frac59,\quad E_2=\frac13.
\]
The core identity is
$2A_2^2-O_2=-2/9=D_2+E_2$.
This shows explicitly how both boundary terms enter the calculation.

\begin{thebibliography}{9}
\raggedright
\bibitem{CampbellLevrie}
John M. Campbell and Paul Levrie,
\emph{On the Basel Problem and the Square of Gregory's Series},
Mathematics Magazine \textbf{97}(2) (2024), 181--186.
\url{https://doi.org/10.1080/0025570X.2024.2312093}.
Used for historical context; the finite proof and bounds are
provided in full above.

\bibitem{ProveItPR}
Vladimir Reshetnikov, \emph{ProveIt: ConciseIPRA},
revision \texttt{b03925cccc76d55b9dc2a00864a5e0544b64aa7f},
\texttt{PR.lean} and \texttt{Deduction.lean}.
\url{https://github.com/VladimirReshetnikov/ProveIt/blob/b03925cccc76d55b9dc2a00864a5e0544b64aa7f/Logic/PrimitiveRecursiveArithmetic/ConciseIPRA/Lean/ConciseIPRA/PR.lean}.

\bibitem{ProveItNative}
Vladimir Reshetnikov, \emph{NativeBiInterpretation.lean},
same revision, theorem \texttt{pra\_iff\_native}.
\url{https://github.com/VladimirReshetnikov/ProveIt/blob/b03925cccc76d55b9dc2a00864a5e0544b64aa7f/Logic/PrimitiveRecursiveArithmetic/ConciseIPRA/Lean/ConciseIPRA/NativeBiInterpretation.lean}.

\bibitem{BaselProblem}
\emph{Basel problem}, user-supplied problem reference.
\url{https://en.wikipedia.org/wiki/Basel_problem}.
The requested theorem is Euler's value $\zeta(2)=\pi^2/6$;
the argument here is self-contained.
\end{thebibliography}
\end{document}
